\documentclass[10pt]{article}

\usepackage[utf8]{inputenc}
\usepackage[T1]{fontenc}
\usepackage{amsmath,amssymb,amsthm}
\usepackage{microtype}
\usepackage[margin=1in]{geometry}

\newtheorem*{theorem}{Theorem}

\title{Partial Fractions and Four Classical Theorems of Number Theory}
\author{Michael D. Hirschhorn}
\date{}

\begin{document}
\maketitle

For some years Jacobi's two-square and four-square theorems have been sources of
considerable fascination to me, so much so that I have continually sought the
simplest, most direct proofs of them. Indeed, I have in the past presented two
proofs of each in \cite{Hirschhorn1982,Hirschhorn1985,Hirschhorn1987,Hirschhorn1999}.
It is my intention to give now what I regard as even simpler proofs of these two
theorems, as well as results due to Dirichlet and Lorenz on representations by a
square plus twice a square and a square plus three times a square, respectively.
All four classical theorems follow from just one partial-fractions expansion
together with special cases of Jacobi's triple-product identity. This work is my
attempt to make rigorous, to simplify, and to present to a wider readership the
work of Lorenz \cite{Lorenz1871}.

With $d_{r,m}(n)$ denoting the number of divisors $d$ of $n$ with
$d\equiv r\pmod m$, the theorems we prove are as follows.

\begin{theorem}[Jacobi, 1828]
The number of representations of the number $n\geq 1$ as the sum of two squares is
\[
4\bigl(d_{1,4}(n)-d_{3,4}(n)\bigr).
\]
\end{theorem}

\begin{theorem}[Dirichlet, 1840]
The number of representations of $n\geq 1$ as the sum of a square and twice a
square is
\[
2\bigl(d_{1,8}(n)+d_{3,8}(n)-d_{5,8}(n)-d_{7,8}(n)\bigr).
\]
\end{theorem}

\begin{theorem}[Lorenz, 1871]
The number of representations of $n\geq 1$ as the sum of a square and three times
a square is
\[
2\bigl(d_{1,3}(n)-d_{2,3}(n)\bigr)
 +4\bigl(d_{4,12}(n)-d_{8,12}(n)\bigr).
\]
\end{theorem}

and

\begin{theorem}[Jacobi, 1829]
The number of representations of $n\geq 1$ as the sum of four squares is
\[
8\sum_{\substack{d\mid n\\4\nmid d}} d.
\]
\end{theorem}

\section*{The proofs}

Before presenting the partial-fractions expansion referred to in the opening
paragraph, let us look at the special cases of Jacobi's triple-product identity
that we require. Jacobi's triple-product identity says that if $a\neq 0$ and
$|q|<1$, then
\[
\prod_{n\geq 1}
 \bigl(1+aq^{2n-1}\bigr)
 \bigl(1+a^{-1}q^{2n-1}\bigr)
 \bigl(1-q^{2n}\bigr)
 =\sum_{n=-\infty}^{\infty}a^nq^{n^2}.
\]
If we set $a=-1$, we find
\[
\begin{aligned}
\sum_{n=-\infty}^{\infty}(-1)^nq^{n^2}
 &=\prod_{n\geq 1}\bigl(1-q^{2n-1}\bigr)^2\bigl(1-q^{2n}\bigr) \\
 &=\prod_{n\geq 1}\frac{\bigl(1-q^n\bigr)^2}{1-q^{2n}}
  =\prod_{n\geq 1}\frac{1-q^n}{1+q^n} \\
 &=\prod_{n\geq 1}
   \frac{1-q^{2n-1}}{1+q^{2n}}
   \frac{1-q^{2n}}{1+q^{2n-1}},
\end{aligned}
\]
and, if we put $-q^2$ for $q$ in this, we find
\[
\begin{aligned}
\sum_{n=-\infty}^{\infty}q^{2n^2}
 &=\prod_{n\geq 1}
   \frac{\bigl(1+q^{4n-2}\bigr)\bigl(1-q^{4n}\bigr)}
        {\bigl(1+q^{4n}\bigr)\bigl(1-q^{4n-2}\bigr)} \\
 &=\prod_{n\geq 1}
   \frac{\bigl(1+q^{4n-2}\bigr)\bigl(1-q^{2n}\bigr)}
        {\bigl(1+q^{4n}\bigr)\bigl(1+q^{2n-1}\bigr)}
   \frac{1+q^{2n}}{1-q^{2n-1}}.
\end{aligned}
\]

Now, it is an easy exercise in partial fractions to show that, provided no term
in a denominator is zero,
\begin{equation}
\frac{1}{2-x}\prod_{n=1}^{N}
 \frac{1-aq^{n-1}x+a^2q^{2n-2}}{1-q^nx+q^{2n}}
 =\frac{A_0}{2-x}
  +\sum_{n=1}^{N}\frac{A_n}{1-q^nx+q^{2n}},
\label{eq:finite-partial-fractions}
\end{equation}
where, with
\[
(a;q)_n=(1-a)(1-aq)\cdots\bigl(1-aq^{n-1}\bigr),
\qquad (a;q)_0=1,
\]
we have
\[
A_0=\frac{(a;q)_N^2}{(q;q)_N^2},
\]
and, for $n\geq 1$,
\[
A_n=
\frac{a^n(a^{-1}q;q)_n(1+q^n)}{(a;q)_n}
\frac{(a;q)_{N+n}(a;q)_{N-n}}
     {(q;q)_{N+n}(q;q)_{N-n}}.
\]
Dividing \eqref{eq:finite-partial-fractions} by $A_0$, assuming
$0<|a|<1$ and $|q|<1$, and letting $N\to\infty$, we obtain
\begin{equation}
\begin{aligned}
&\frac{1}{2-x}\prod_{n\geq 1}
 \frac{1-aq^{n-1}x+a^2q^{2n-2}}{1-q^nx+q^{2n}}
 \frac{(1-q^n)^2}{(1-aq^{n-1})^2} \\
&\qquad=\frac{1}{2-x}
 +\sum_{n\geq 1}
  \frac{a^n(a^{-1}q;q)_n(1+q^n)}
       {(a;q)_n(1-q^nx+q^{2n})},
\end{aligned}
\label{eq:infinite-partial-fractions}
\end{equation}
again provided no term in a denominator is zero.

If in \eqref{eq:infinite-partial-fractions} we set $a=-q$ and $x=0$, we find
\[
\prod_{n\geq 1}\frac{(1-q^n)^2}{(1+q^n)^2}
 =1+4\sum_{n\geq 1}\frac{(-1)^nq^n}{1+q^{2n}},
\]
or
\begin{equation}
\left(\sum_{n=-\infty}^{\infty}(-1)^nq^{n^2}\right)^2
 =1+4\sum_{n\geq 1}\frac{(-1)^nq^n}{1+q^{2n}}.
\label{eq:two-squares-base}
\end{equation}
If in \eqref{eq:two-squares-base} we put $-q$ for $q$, we obtain
\[
\begin{aligned}
\left(\sum_{n=-\infty}^{\infty}q^{n^2}\right)^2
 &=1+4\sum_{n\geq 1}\frac{q^n}{1+q^{2n}} \\
 &=1+4\sum_{n\geq 1}
   \left(\frac{q^n}{1-q^{4n}}-\frac{q^{3n}}{1-q^{4n}}\right) \\
 &=1+4\sum_{\substack{k\geq 1\\\ell\geq 0}}
   \left(q^{k(4\ell+1)}-q^{k(4\ell+3)}\right) \\
 &=1+4\sum_{n\geq 1}\bigl(d_{1,4}(n)-d_{3,4}(n)\bigr)q^n,
\end{aligned}
\]
from which Jacobi's two-square theorem follows.

If in \eqref{eq:infinite-partial-fractions} we set $a=-q$ and $x=-2$, we find
\[
\prod_{n\geq 1}\frac{(1-q^n)^4}{(1+q^n)^4}
 =1+8\sum_{n\geq 1}\frac{(-1)^nq^n}{(1+q^n)^2},
\]
or
\begin{equation}
\left(\sum_{n=-\infty}^{\infty}(-1)^nq^{n^2}\right)^4
 =1+8\sum_{n\geq 1}\frac{(-1)^nq^n}{(1+q^n)^2}.
\label{eq:four-squares-base}
\end{equation}
If in \eqref{eq:four-squares-base} we put $-q$ for $q$, we obtain
\[
\begin{aligned}
\left(\sum_{n=-\infty}^{\infty}q^{n^2}\right)^4
 &=1+8\sum_{\substack{n\geq 1\\n\text{ even}}}
       \frac{q^n}{(1+q^n)^2}
   +8\sum_{\substack{n\geq 1\\n\text{ odd}}}
       \frac{q^n}{(1-q^n)^2} \\
 &=1+8\sum_{n\geq 1}\frac{q^n}{(1-q^n)^2}
   -8\sum_{\substack{n\geq 1\\n\text{ even}}}
    \left(\frac{q^n}{(1-q^n)^2}-\frac{q^n}{(1+q^n)^2}\right) \\
 &=1+8\sum_{n\geq 1}\frac{q^n}{(1-q^n)^2}
   -32\sum_{n\geq 1}\frac{q^{4n}}{(1-q^{4n})^2} \\
 &=1+8\sum_{n\geq 1}
   \left(\sigma(n)-4\sigma\!\left(\frac{n}{4}\right)\right)q^n \\
 &=1+8\sum_{n\geq 1}
   \left(\sum_{\substack{d\mid n\\4\nmid d}}d\right)q^n,
\end{aligned}
\]
where $\sigma(m)$ is the sum of the positive divisors of $m$, with
$\sigma(x)=0$ when $x$ is not a positive integer. This proves Jacobi's
four-square theorem.

If in \eqref{eq:infinite-partial-fractions} we set $a=-q$ and $x=1$, we find
\[
\prod_{n\geq 1}
 \frac{1-q^n}{1+q^n}\frac{1-q^{3n}}{1+q^{3n}}
 =1+2\sum_{n\geq 1}\frac{(-1)^nq^n}{1-q^n+q^{2n}},
\]
or
\begin{equation}
\left(\sum_{n=-\infty}^{\infty}(-1)^nq^{n^2}\right)
\left(\sum_{n=-\infty}^{\infty}(-1)^nq^{3n^2}\right)
 =1+2\sum_{n\geq 1}\frac{(-1)^nq^n}{1-q^n+q^{2n}}.
\label{eq:lorenz-base}
\end{equation}
If in \eqref{eq:lorenz-base} we put $-q$ for $q$, we obtain
\[
\begin{aligned}
&\left(\sum_{n=-\infty}^{\infty}q^{n^2}\right)
 \left(\sum_{n=-\infty}^{\infty}q^{3n^2}\right) \\
&\quad=1+2\sum_{n\geq 1}
 \frac{q^n}{1-(-1)^nq^n+q^{2n}} \\
&\quad=1+2\sum_{\substack{n\geq 1\\n\text{ even}}}
 \frac{q^n(1+q^n)}{1+q^{3n}}
 +2\sum_{\substack{n\geq 1\\n\text{ odd}}}
 \frac{q^n(1-q^n)}{1-q^{3n}} \\
&\quad=1+2\sum_{n\geq 1}\frac{q^n(1-q^n)}{1-q^{3n}} \\
&\qquad\quad
 +2\sum_{\substack{n\geq 1\\n\text{ even}}}
 \left(\frac{q^n(1+q^n)}{1+q^{3n}}
       -\frac{q^n(1-q^n)}{1-q^{3n}}\right) \\
&\quad=1+2\sum_{n\geq 1}\frac{q^n(1-q^n)}{1-q^{3n}}
 +4\sum_{\substack{n\geq 1\\n\text{ even}}}
 \frac{q^{2n}(1-q^{2n})}{1-q^{6n}} \\
&\quad=1+2\sum_{n\geq 1}\frac{q^n(1-q^n)}{1-q^{3n}}
 +4\sum_{n\geq 1}\frac{q^{4n}(1-q^{4n})}{1-q^{12n}} \\
&\quad=1+2\sum_{n\geq 1}\bigl(d_{1,3}(n)-d_{2,3}(n)\bigr)q^n \\
&\qquad\quad
 +4\sum_{n\geq 1}\bigl(d_{4,12}(n)-d_{8,12}(n)\bigr)q^n,
\end{aligned}
\]
from which Lorenz's theorem follows.

If in \eqref{eq:infinite-partial-fractions} we put $q^2$ for $q$, set $a=-q$,
and set $x=0$, we find
\[
\prod_{n\geq 1}
 \frac{1+q^{4n-2}}{1+q^{4n}}
 \frac{(1-q^{2n})^2}{(1+q^{2n-1})^2}
 =1+2\sum_{n\geq 1}
 \frac{(-1)^nq^n(1+q^{2n})}{1+q^{4n}},
\]
or
\begin{equation}
\left(\sum_{n=-\infty}^{\infty}q^{2n^2}\right)
\left(\sum_{n=-\infty}^{\infty}(-1)^nq^{n^2}\right)
 =1+2\sum_{n\geq 1}
 \frac{(-1)^nq^n(1+q^{2n})}{1+q^{4n}}.
\label{eq:dirichlet-base}
\end{equation}
If in \eqref{eq:dirichlet-base} we put $-q$ for $q$, we obtain
\[
\begin{aligned}
\left(\sum_{n=-\infty}^{\infty}q^{n^2}\right)
\left(\sum_{n=-\infty}^{\infty}q^{2n^2}\right)
 &=1+2\sum_{n\geq 1}\frac{q^n(1+q^{2n})}{1+q^{4n}} \\
 &=1+2\sum_{n\geq 1}
 \left(\frac{q^n}{1-q^{8n}}
      +\frac{q^{3n}}{1-q^{8n}}
      -\frac{q^{5n}}{1-q^{8n}}
      -\frac{q^{7n}}{1-q^{8n}}\right) \\
 &=1+2\sum_{n\geq 1}
 \bigl(d_{1,8}(n)+d_{3,8}(n)-d_{5,8}(n)-d_{7,8}(n)\bigr)q^n,
\end{aligned}
\]
from which Dirichlet's theorem follows.

\begin{thebibliography}{9}

\bibitem{Hirschhorn1982}
M.~D. Hirschhorn,
``A simple proof of Jacobi's four-square theorem,''
\emph{J. Austral. Math. Soc. Ser. A} \textbf{32} (1982), 61--67.

\bibitem{Hirschhorn1985}
M.~D. Hirschhorn,
``A simple proof of Jacobi's two-square theorem,''
\emph{Amer. Math. Monthly} \textbf{92} (1985), 579--580.

\bibitem{Hirschhorn1987}
M.~D. Hirschhorn,
``A simple proof of Jacobi's four-square theorem,''
\emph{Proc. Amer. Math. Soc.} \textbf{101} (1987), 436--438.

\bibitem{Hirschhorn1999}
M.~D. Hirschhorn,
``Jacobi's two-square theorem and related identities,''
\emph{Ramanujan J.} \textbf{3} (1999), 153--158.

\bibitem{Lorenz1871}
L.~Lorenz,
``Contribution à la théorie des nombres,''
in \emph{\OE uvres scientifiques de L. Lorenz},
revues et annotées par H.~Valentiner, vol.~II,
Lehmann \& Stage, Copenhague, 1904, pp.~403--431.

\end{thebibliography}

\bigskip
\noindent\textit{School of Mathematics, UNSW, Sydney 2052, Australia}\\
\texttt{m.hirschhorn@unsw.edu.au}

\end{document}
